% Independent single-column Figure 2; native size 3.34 x 2.35 inches.
% CPU-only static estimates; see reuse-aiconfigurator-max-context-h100/.
\begin{figure}[t]
  \centering
  \includegraphics[width=\linewidth]{figures/motivation-row/motivation-recompute.pdf}
  \caption{Recovery estimates up to each model's context limit.
    Color denotes path; line style and marker denote model.
    Qwen14B, Llama70B, Qwen235B (Instruct-2507), and DeepSeek-V3.2
    use 1/4/16/32 H100 GPUs and end at 128K (YaRN)/128K/1,010,000/160K,
    respectively. BF16, batch one; total context includes the known
    64-token tail. K~=~1,024; M~=~$10^6$.
    Recompute uses 8K chunks; DRAM restores full state at an assumed
    64\,GB/s per GPU before tail compute, without overlap.
    All values are static estimates: Qwen beyond 256K projects full
    attention without DCA acceleration; DSA uses analytical prefix
    correction and table extrapolation beyond 32K.
    Detection, queueing, and software takeover are excluded.}
  \label{fig:motivation-recompute}
  \Description{One plot compares twelve estimated recovery curves on
    logarithmic context and latency axes. Blue, orange, and red encode
    HBM reuse, DRAM reuse, and Recompute. Circles, squares, triangles,
    and diamonds encode Qwen3-14B, Llama-3.1-70B, Qwen3-235B-A22B-Instruct-2507,
    and DeepSeek-V3.2, respectively. The latter two are MoE models;
    DeepSeek additionally uses sparse attention. Curves stop at each
    model's configured context limit. At 1,010,000 total tokens,
    Qwen 2507 costs an estimated 70 milliseconds for HBM reuse,
    830 milliseconds for DRAM reuse, and 220 seconds for chunked
    full-attention recomputation. GPU allocations differ across models.}
\end{figure}
